%%%%%%%%%%%%%%%%%%%%%%%%%%%%%%%%%%%%%%%%%%%%%%%%%%%%%%%%%%%%%%%%%%%
%                                                                 %
%                          CHAPTER SIX                            %
%                                                                 %
%%%%%%%%%%%%%%%%%%%%%%%%%%%%%%%%%%%%%%%%%%%%%%%%%%%%%%%%%%%%%%%%%%%

\chapter{Συμπεράσματα}
\label{ch:conclusions}

Η πειραματική αξιολόγηση του Κεφαλαίου~\ref{ch:experiments} ανέδειξε τέσσερα κεντρικά
συμπεράσματα. Πρώτον, η αναδρομική αρχιτεκτονική LSTM υπερτερεί σαφώς της απλής MLP εκτός
δείγματος και, σε αντίθεση με τις έτοιμες βιβλιοθήκες ενισχυτικής μάθησης, η προσαρμοσμένη
υλοποίηση δεν εμφανίζει την παθολογία ακραίας υπερπροσαρμογής. Παράλληλα, αναδείχθηκε η
παθητική κατάρρευση ως πραγματικό και επαναλαμβανόμενο φαινόμενο, το οποίο αντιμετωπίζεται
αποτελεσματικά μόνο με την ανταμοιβή ενεργής απόδοσης. Αφαιρώντας το beta της αγοράς και
επιβάλλοντας σχετική επιλογή νικητών, η ανταμοιβή αυτή επέτυχε λόγο Sharpe 0,767 και
απόδοση $+33{,}5\%$, με χαμηλότερη διασπορά από τις εναλλακτικές διαμορφώσεις.

Δεύτερον, η επιλογή της διατύπωσης των περιορισμών ρευστότητας αποδείχθηκε καθοριστική.
Ακατάλληλες διατυπώσεις μπορούν να παραγάγουν γωνιακές λύσεις αδρανούς κράτησης μετρητών,
όπως η διατύπωση EMA με φαινομενικά υψηλό Sharpe 1,250, οι οποίες αποκαλύπτονται ως
κατάρρευση μόνο μέσω διαγνωστικής ανάλυσης συμπεριφοράς και όχι μέσω των τυπικών μετρικών.
Το εύρημα αυτό υπογραμμίζει ότι η αξιολόγηση ενός πράκτορα συναλλαγών δεν μπορεί να
στηριχθεί αποκλειστικά σε αριθμητικούς δείκτες απόδοσης. Η soft-start διατύπωση, σε
συνδυασμό με βελτιστοποίηση μέσω CPCV και Optuna, απέφυγε αυτές τις παθολογίες και βελτίωσε
ουσιαστικά το προφίλ κινδύνου-απόδοσης, φτάνοντας σε Sharpe 1,072 και απόδοση $+50{,}3\%$.

Τρίτον, επιβεβαιώθηκε η βασική υπόθεση της εργασίας ως προς την αξία του κειμενικού σήματος,
αλλά υπό συγκεκριμένη συνθήκη. Το σήμα συναισθήματος FinGPT προσφέρει σημαντική αξία κυρίως
μέσω της αλληλεπίδρασής του με τον Λαγκρανζιανό περιορισμό liquify, με βελτίωση $+0{,}759$
στον λόγο Sharpe έναντι της ίδιας διαμόρφωσης χωρίς συναίσθημα, ενώ ως αυτόνομο σήμα αποφέρει
οριακή μόνο βελτίωση $+0{,}030$. Η συνέργεια ερμηνεύεται ως συνδυασμός της πίεσης για επένδυση
που επιβάλλει ο περιορισμός και της κατεύθυνσης επένδυσης που παρέχει το συναίσθημα. Επιπλέον,
η σύγκριση των δύο τρόπων ενσωμάτωσης έδειξε ότι η αξία του σήματος εντοπίζεται στο
πληροφοριακό του περιεχόμενο εντός του χώρου καταστάσεων και όχι στο κανάλι προσαρμογής των
ενεργειών, καθώς οι δύο προσεγγίσεις απέδωσαν στατιστικά ισοδύναμα αποτελέσματα.

Τέταρτον, η βελτιστοποίηση μέσω CPCV και Optuna αποδείχθηκε αναγκαία για την ανάδειξη της
πλήρους αξίας κάθε διαμόρφωσης, εξασφαλίζοντας ταυτόχρονα ότι η αναζήτηση δεν οδηγεί σε
υπερπροσαρμογή στην περίοδο εκπαίδευσης. Συνολικά, ο βέλτιστος πράκτορας της εργασίας,
δηλαδή ο Lagrangian liquify 15\%/70\% με χαρακτηριστικά FinGPT one-hot βελτιστοποιημένος μέσω
CPCV, επιτυγχάνει λόγο Sharpe 1,485 και απόδοση $+47{,}9\%$ από αρχικό κεφάλαιο
$1.000.000\$$, υπερβαίνοντας κατά πολύ τόσο την ισόποση αγορά-και-κράτηση (Sharpe 0,267,
απόδοση $-7{,}8\%$) όσο και τον πράκτορα ενεργής απόδοσης χωρίς περιορισμούς ή συναίσθημα
(Sharpe 0,767, απόδοση $+33{,}5\%$). Το αποτέλεσμα αυτό προκύπτει από τη συνδυαστική εφαρμογή
και των τριών αξόνων της εργασίας, της αρχιτεκτονικής και της ανταμοιβής, του Λαγκρανζιανού
περιορισμού ρευστότητας και του σήματος συναισθήματος, και δεν μπορεί να αποδοθεί σε κανέναν
από αυτούς μεμονωμένα.

Τα ευρήματα της εργασίας ανοίγουν αρκετές κατευθύνσεις για μελλοντική διερεύνηση, οι οποίες
προκύπτουν τόσο από τις επιτυχίες όσο και από τους περιορισμούς της παρούσας προσέγγισης.

\textbf{Ευρωστία σε πολλαπλά καθεστώτα αγοράς.} Η αξιολόγηση περιορίστηκε σε μία
περίοδο εκτός δείγματος (2024--2025). Η εφαρμογή κυλιόμενης επαναξιολόγησης (walk-forward)
σε περισσότερες και μακρύτερες περιόδους, που να καλύπτουν διακριτά καθεστώτα ανόδου,
πτώσης και πλάγιας κίνησης, θα ενίσχυε τη γενίκευση των συμπερασμάτων.

\textbf{Βελτιωμένες διατυπώσεις περιορισμών.} Ορισμένες διατυπώσεις, όπως οι
διπλόπλευρες ζώνες μετρητών, κατέρρευσαν λόγω αντικρουόμενων κινήτρων. Η αυτόματη και
προσαρμοστική ρύθμιση των ορίων του περιορισμού, καθώς και εναλλακτικές διατυπώσεις CMDP,
αποτελούν φυσική επέκταση για την αποφυγή τέτοιων παθολογιών.

\textbf{Πλουσιότερο σήμα συναισθήματος.} Το σήμα χρησιμοποιήθηκε σε μορφή
συγκεντρωτικής κωδικοποίησης one-hot. Η αξιοποίηση συνεχών σκορ, υψηλότερης χρονικής
ανάλυσης, ανά γεγονός πληροφορίας, καθώς και η σύγκριση εναλλακτικών γλωσσικών μοντέλων,
θα μπορούσε να αναδείξει πρόσθετη πληροφοριακή αξία.

\textbf{Μηχανισμοί σύζευξης συναισθήματος και ενεργειών.} Δεδομένου ότι η σύζευξη
με τον χώρο ενεργειών δεν προσέφερε στατιστικά σημαντικό πλεονέκτημα έναντι της αμιγούς
χρήσης ως χαρακτηριστικού κατάστασης, αξίζει η διερεύνηση πιο εκφραστικών μηχανισμών
σύζευξης που να αξιοποιούν το σήμα πέρα από την απλή κανονικοποίηση των ενεργειών.

\textbf{Ρεαλισμός εκτέλεσης.} Η ενσωμάτωση μοντέλων ολίσθησης τιμής (slippage) και
επίδρασης στην αγορά (market impact), η επέκταση σε ευρύτερο σύνολο περιουσιακών στοιχείων
και η αξιολόγηση σε συνθήκες προσομοιωμένης ζωντανής συναλλαγής θα προσέγγιζαν περισσότερο
τις πραγματικές συνθήκες λειτουργίας.

\textbf{Στατιστική ισχύς και επέκταση πεδίου.} Η αύξηση του αριθμού των παραγόντων
τυχαιοποίησης και η εφαρμογή τυπικών ελέγχων στατιστικής σημαντικότητας, λόγω της υψηλής
διασποράς των αποτελεσμάτων, καθώς και η μεταφορά της μεθοδολογίας σε άλλες κατηγορίες
περιουσιακών στοιχείων, όπως μετοχές και συνάλλαγμα, θα ενίσχυαν την αξιοπιστία και την
ευρύτητα των ευρημάτων.

Συνοψίζοντας, η εργασία κατέδειξε ότι ο συνδυασμός μιας αναδρομικής αρχιτεκτονικής, μιας
ανταμοιβής που αποτρέπει την παθητική κατάρρευση, ενός κατάλληλα διατυπωμένου περιορισμού
ρευστότητας και ενός σήματος συναισθήματος από εξειδικευμένο γλωσσικό μοντέλο μπορεί να
οδηγήσει σε πράκτορα συναλλαγών με σαφώς ανώτερη προσαρμοσμένη στον κίνδυνο απόδοση. Το
κρισιμότερο μεθοδολογικό δίδαγμα είναι ότι η αξία κάθε επιμέρους συνιστώσας αναδεικνύεται
μόνο μέσω της αλληλεπίδρασής της με τις υπόλοιπες και μόνο υπό αυστηρή αξιολόγηση εκτός
δείγματος που συμπληρώνεται από διαγνωστική ανάλυση συμπεριφοράς.
